\chapter{Formas hermíticas}\label{ch:b2:hermitian}

Los espacios vectoriales complejos tienen su propia geometría de
producto escalar, con un giro: linealidad en una variable y
\emph{semi}linealidad en la otra. La recompensa por aceptar ese giro es
una teoría espectral aún más limpia que la real: los
\omterm{def:b2:hermitian:adjoint}{endomorfismos hermíticos} tienen \omterm{def:b2:reduction:eigen}{valores propios} reales, los
\omterm{def:b2:hermitian:adjoint}{unitarios} los tienen de módulo $1$, y ambos se diagonalizan en bases
ortonormales. Este capítulo breve reproduce sobre $\C$ el programa
euclídeo del~\cref{ch:b2:quadratic}.

\section{Productos escalares hermíticos}

\begin{definition}\label{def:b2:hermitian:def}
Un \emph{producto escalar hermítico}\index{producto escalar
hermítico} sobre un espacio vectorial complejo $E$ es una aplicación
$\langle\cdot,\cdot\rangle \colon E \times E \to \C$ que es lineal en
la segunda variable, \omterm{def:b2:quadratic:adjoint}{simétrica} conjugada ($\langle y, x\rangle =
\conj{\langle x, y\rangle}$, luego semilineal en la primera
variable) y definida positiva ($\langle x, x\rangle > 0$ para $x \neq
0$). El ejemplo estándar sobre $\C^n$:
\[
\langle x, y \rangle = \sum_{i=1}^{n} \conj{x_i}\, y_i ;
\]
sobre funciones \omterm{def:b2:metric:continuity}{continuas}, $\langle f, g\rangle = \int_a^b \conj
f\,g$. \omterm{def:b2:nvs:norm}{Norma}: $\norm x = \sqrt{\langle x,x\rangle}$; un espacio
complejo de dimensión finita así equipado es un \emph{espacio
hermítico}\index{espacio hermítico}.
\end{definition}

\begin{example}[Primeros cálculos]\label{ex:b2:hermitian:firstcomp}
En $\C^2$, tomemos $x = (1+\iu,\ 2-\iu)$ e $y = (\iu,\ 1)$.
Entonces
\[
\norm x^2 = \abs{1+\iu}^2 + \abs{2-\iu}^2 = 2 + 5 = 7,
\qquad
\norm y^2 = 1 + 1 = 2,
\]
y, conjugando el primer argumento,
\[
\langle x, y\rangle
= \conj{(1+\iu)}\,\iu + \conj{(2-\iu)}\cdot1
= (1-\iu)\iu + (2+\iu)
= (\iu + 1) + (2 + \iu) = 3 + 2\iu .
\]
Cauchy--Schwarz se verifica: $\abs{\langle x, y\rangle}^2 = 9 + 4 =
13 \leq 14 = \norm x^2\norm y^2$; cerca de la igualdad, porque $x$
está cerca de ser un múltiplo de $y$. Obsérvese también que $\langle
y, x\rangle = \conj{3 + 2\iu} = 3 - 2\iu$: la simetría conjugada en
acción, y la razón de que $\langle x, x\rangle$ sea siempre real.
\end{example}

\begin{theorem}[Cauchy--Schwarz, caso complejo]\label{thm:b2:hermitian:cs}
$\abs{\langle x, y\rangle} \leq \norm x \norm y$, con igualdad si y
solo si $x, y$ son linealmente dependientes; y $\norm\cdot$ es una
\omterm{def:b2:nvs:norm}{norma}. Además existen bases ortonormales (Gram--Schmidt funciona
palabra por palabra), con
\[
x = \sum_i \langle e_i, x\rangle\, e_i,
\qquad
\norm x^2 = \sum_i \abs{\langle e_i, x\rangle}^2 .
\]
\end{theorem}

\begin{proof}
Para $y \neq 0$ y $t \in \C$: $0 \leq \norm{x - ty}^2 = \norm x^2 -
2\Re\bigl(\conj t\langle y, x\rangle\bigr) + \abs t^2\norm y^2$.
Elíjase $t = \frac{\langle y, x\rangle}{\norm y^2}$:
\[
0 \leq \norm x^2 - \frac{\abs{\langle y, x\rangle}^2}{\norm y^2},
\]
que es la desigualdad; la igualdad fuerza $x = ty$. Desigualdad
triangular, en detalle:
\[
\norm{x + y}^2 = \norm x^2 + 2\,\Re\langle x, y\rangle
+ \norm y^2
\leq \norm x^2 + 2\,\abs{\langle x, y\rangle} + \norm y^2
\leq \bigl(\norm x + \norm y\bigr)^2 ,
\]
usando $\Re z \leq \abs z$ y luego Cauchy--Schwarz; la homogeneidad y
la separación son inmediatas, de modo que $\norm\cdot$ es una
\omterm{def:b2:nvs:norm}{norma}. Gram--Schmidt: como en el caso real, con las conjugaciones
colocadas según la definición (nótese el convenio: nuestros productos
son semilineales en la \emph{primera} entrada, así que las
coordenadas son $\langle e_i, x\rangle$; el ejemplo siguiente ejecuta
el algoritmo una vez por completo).
\end{proof}

\begin{example}[Gram--Schmidt complejo, ejecutado por completo]\label{ex:b2:hermitian:gramschmidt}
Ortonormalicemos la base $v_1 = (1, \iu)$, $v_2 = (0, 1)$ de $\C^2$.
Primer vector: $\norm{v_1}^2 = \abs1^2 + \abs\iu^2 = 2$, luego $e_1 =
\frac{1}{\sqrt2}(1, \iu)$. Proyectemos $v_2$, con la conjugación en
la primera entrada:
\[
\langle e_1, v_2\rangle
= \frac{1}{\sqrt2}\bigl(\conj{1}\cdot0 +
\conj{\iu}\cdot1\bigr)
= \frac{-\iu}{\sqrt2} ,
\qquad
v_2 - \langle e_1, v_2\rangle e_1
= (0,1) + \frac{\iu}{2}\,(1, \iu)
= \Bigl(\frac\iu2,\ \frac12\Bigr) .
\]
Su \omterm{def:b2:nvs:norm}{norma} es $\sqrt{\frac14 + \frac14} = \frac{1}{\sqrt2}$: $e_2 =
\frac{1}{\sqrt2}(\iu, 1)$. Comprobación: $\langle e_1, e_2\rangle =
\frac12(\conj1\cdot\iu + \conj\iu\cdot1) = \frac12(\iu - \iu) = 0$.
Coordenadas de $v_2$ en la nueva base: $v_2 = \langle e_1,
v_2\rangle e_1 + \langle e_2, v_2\rangle e_2$ con $\langle e_2,
v_2\rangle = \frac{1}{\sqrt2}$; atención al orden: $\langle v_2,
e_1\rangle$ daría el coeficiente \emph{conjugado}. Moraleja: el
algoritmo es el euclídeo palabra por palabra; la única trampa está en
dónde cae la conjugación, y calcular $\norm{v_2 - \text{proy}}^2 > 0$
usa en silencio la positividad, el axioma que hace funcionar toda la
geometría.
\end{example}

\section{Adjunto, endomorfismos hermíticos y unitarios}

\begin{definition}\label{def:b2:hermitian:adjoint}
El \emph{adjunto} $u^*$ de $u \in \mathcal{L}(E)$ se define por
$\langle u^*(x), y\rangle = \langle x, u(y)\rangle$; en una base
ortonormal, $\operatorname{Mat}(u^*) = \conj{A}^{\mathsf T} =:
A^{\dagger}$ (la \omterm{def:b2:linalg:transpose}{traspuesta} conjugada). En efecto, si $B = (b_{ij})$
es la matriz de $u^*$ en la base ortonormal $(e_i)$, entonces
$b_{ij} = \langle e_i, u^*(e_j)\rangle$, y la identidad que lo define
da
\[
\conj{b_{ij}}
= \langle u^*(e_j), e_i\rangle
= \langle e_j, u(e_i)\rangle = a_{ji} ,
\qquad\text{luego}\qquad
B = \conj{A}^{\mathsf T} .
\] Se dice que $u$ es
\emph{hermítico}\index{endomorfismo hermítico} cuando $u^* = u$
($A^\dagger = A$), \emph{unitario}\index{grupo unitario} cuando $u^*u
= \mathrm{id}$ ($A^\dagger A = I$: el grupo $U(n)$) y \emph{normal}
cuando $u^*u = uu^*$.
\end{definition}

\begin{proposition}\label{prop:b2:hermitian:eigenvalues}
Los \omterm{def:b2:reduction:eigen}{valores propios} de un \omterm{def:b2:hermitian:adjoint}{endomorfismo hermítico} son \emph{reales};
los \omterm{def:b2:reduction:eigen}{valores propios} de un endomorfismo \omterm{def:b2:hermitian:adjoint}{unitario} tienen
\emph{módulo $1$}; y en ambos casos los \omterm{def:b2:reduction:eigen}{subespacios propios} de
\omterm{def:b2:reduction:eigen}{valores propios} distintos son ortogonales.
\end{proposition}

\begin{proof}
Caso \omterm{def:b2:hermitian:adjoint}{hermítico}, con $u(x) = \lambda x$, $x \neq 0$:
\[
\lambda \norm x^2 = \langle x, u(x)\rangle = \langle u(x), x\rangle
= \conj{\langle x, u(x)\rangle} = \conj\lambda\,\norm x^2 ,
\]
luego $\lambda \in \R$. Caso \omterm{def:b2:hermitian:adjoint}{unitario}: $\norm{u(x)} = \norm x$ (por
$u^*u = \mathrm{id}$), de modo que $\abs\lambda\norm x = \norm x$.
Ortogonalidad (caso \omterm{def:b2:hermitian:adjoint}{hermítico}): $\lambda\langle x, y\rangle = \langle
u(x), y\rangle = \langle x, u(y)\rangle = \mu\langle x, y\rangle$
para \omterm{def:b2:reduction:eigen}{vectores propios} con $\lambda \neq \mu$ reales. Caso
\omterm{def:b2:hermitian:adjoint}{unitario}, en detalle: para $u(x) = \lambda x$ y $u(y) = \mu y$ con
$\lambda \neq \mu$ (ambos de módulo $1$),
\[
\langle x, y\rangle = \langle u(x), u(y)\rangle
= \conj\lambda\mu\,\langle x, y\rangle ,
\]
y $\conj\lambda\mu = \frac{\mu}{\lambda} \neq 1$: el factor no es
$1$, luego $\langle x, y\rangle = 0$.
\end{proof}

\begin{example}[Una matriz antihermítica, diagonalizada]\label{ex:b2:hermitian:skewexample}
$A = \begin{pmatrix} 0 & -2\\ 2 & 0\end{pmatrix}$ cumple $A^\dagger =
A^{\mathsf T} = -A$: es antihermítica (y también real antisimétrica;
sobre $\R$ no tiene ningún \omterm{def:b2:reduction:eigen}{valor propio}). \omterm{def:b2:reduction:charpoly}{Polinomio característico}
$X^2 + 4$: \omterm{def:b2:reduction:eigen}{valores propios} $\pm2\iu$, imaginarios puros, como predice
en general el~\cref{exo:b2:hermitian:9}. \omterm{def:b2:reduction:eigen}{Vectores propios}: $(A - 2\iu
I)v = 0$ da $v_1 = \frac{1}{\sqrt2}(1, \iu)$, y $v_2 =
\frac{1}{\sqrt2}(1, -\iu)$ para $-2\iu$; son ortogonales:
\[
\langle v_1, v_2\rangle
= \tfrac12\bigl(\conj{1}\cdot1 +
\conj{\iu}\cdot(-\iu)\bigr)
= \tfrac12(1 - 1) = 0 .
\]
Así pues, $A = U\operatorname{diag}(2\iu, -2\iu)\,U^\dagger$ con $U =
(v_1\ v_2)$ \omterm{def:b2:hermitian:adjoint}{unitaria}. Moraleja: $H = -\iu A = \begin{pmatrix} 0 &
2\iu\\ -2\iu & 0\end{pmatrix}$ es \omterm{def:b2:hermitian:adjoint}{hermítica} con \omterm{def:b2:reduction:eigen}{espectro} real
$\{\pm2\}$ y los \emph{mismos} \omterm{def:b2:reduction:eigen}{vectores propios}; es la biyección $u
\mapsto \iu u$ entre \omterm{def:b2:hermitian:adjoint}{endomorfismos hermíticos} y antihermíticos
(\cref{exo:b2:hermitian:9}), vista matriz a matriz. Sobre $\R$, esa
misma $A$ es una rotación con homotecia sin ningún \omterm{def:b2:reduction:eigen}{vector propio}, y
solo el paso a $\C$ revela su forma normal.
\end{example}

\begin{theorem}[Teorema espectral hermítico]\label{thm:b2:hermitian:spectral}
Todo \omterm{def:b2:hermitian:adjoint}{endomorfismo hermítico} de un \omterm{def:b2:hermitian:def}{espacio hermítico} tiene una base
ortonormal de \omterm{def:b2:reduction:eigen}{vectores propios} (con \omterm{def:b2:reduction:eigen}{valores propios} reales):
$A^\dagger = A$ implica $A = U D U^{\dagger}$ con $U \in U(n)$ y $D$
diagonal real.\index{teorema espectral!hermítico}
\end{theorem}

\begin{proof}
Sobre $\C$, el \omterm{def:b2:reduction:charpoly}{polinomio característico} se escinde: existe un
\omterm{def:b2:reduction:eigen}{vector propio} $e_1$ (\cref{ch:b2:reduction}), sin necesidad de ningún
argumento de \omterm{def:b2:metric:compact}{compacidad}, una ventaja de $\C$. Normalicémoslo. Su
complemento ortogonal $F = e_1^\perp$ es estable: para $x \perp e_1$,
\[
\langle e_1, u(x)\rangle = \langle u(e_1), x\rangle
= \lambda_1\langle e_1, x\rangle = 0
\]
($\lambda_1$ real). La restricción es \omterm{def:b2:hermitian:adjoint}{hermítica}; hágase inducción
sobre la dimensión y concaténese. En detalle: la restricción $u|_F$
es un endomorfismo del \omterm{def:b2:hermitian:def}{espacio hermítico} $F$ (de dimensión $n - 1$)
con $\langle u|_F(x), y\rangle = \langle x, u|_F(y)\rangle$ heredado
de $u$; la hipótesis de inducción proporciona una base ortonormal
$(e_2, \dots, e_n)$ de $F$ formada por \omterm{def:b2:reduction:eigen}{vectores propios}, y $(e_1,
e_2, \dots, e_n)$ es ortonormal en $E$ ($e_1 \perp F$) y está
formada por \omterm{def:b2:reduction:eigen}{vectores propios} de $u$. Traducción matricial: las
columnas de $U$ son los $e_i$, $U^\dagger U = I$ expresa su
ortonormalidad y $AU = UD$ recoge las ecuaciones de los
\omterm{def:b2:reduction:eigen}{valores propios}, de donde $A = UDU^\dagger$ con $D$ diagonal real
(\cref{prop:b2:hermitian:eigenvalues}).
\end{proof}

\begin{example}\label{ex:b2:hermitian:pauli}
$A = \begin{pmatrix} 0 & -\iu\\ \iu & 0\end{pmatrix}$ es \omterm{def:b2:hermitian:adjoint}{hermítica}
($A^\dagger = A$): sus \omterm{def:b2:reduction:eigen}{valores propios} salen de $\chi_A = X^2 - 1$:
$\pm 1$ (reales, como se prometió), con \omterm{def:b2:reduction:eigen}{vectores propios}
ortonormales $\frac{1}{\sqrt2}(1, \iu)^{\mathsf T}$ y
$\frac{1}{\sqrt2}(1, -\iu)^{\mathsf T}$. (Los físicos conocen $A$
como matriz de Pauli; que los \omterm{def:b2:reduction:eigen}{espectros} \omterm{def:b2:hermitian:adjoint}{hermíticos} sean reales es la
razón de que los observables cuánticos se modelen mediante operadores
\omterm{def:b2:hermitian:adjoint}{hermíticos}.)
\end{example}

\begin{example}[Una matriz hermítica definida positiva, resuelta]\label{ex:b2:hermitian:pdrun}
$A = \begin{pmatrix} 2 & 1-\iu\\ 1+\iu & 3\end{pmatrix}$: es
\omterm{def:b2:hermitian:adjoint}{hermítica}, pues la diagonal es real y las entradas fuera de ella son
conjugadas. \omterm{def:b2:reduction:charpoly}{Polinomio característico}:
\[
(2-\lambda)(3-\lambda) - \abs{1-\iu}^2
= \lambda^2 - 5\lambda + 4
= (\lambda - 1)(\lambda - 4) :
\]
\omterm{def:b2:reduction:eigen}{espectro} $\{1, 4\}$, real y positivo; $A$ es definida positiva.
\omterm{def:b2:reduction:eigen}{Vectores propios}: para $\lambda = 4$, el sistema $(A - 4I)v = 0$ da
$v_4 = (1 - \iu,\ 2)$ (compruébese la segunda fila: $(1+\iu)(1-\iu) -
2 = 0$); para $\lambda = 1$, $v_1 = (1 - \iu,\ -1)$. Ortogonalidad,
con la conjugación en la primera entrada:
\[
\langle v_4, v_1\rangle
= \conj{(1-\iu)}\,(1-\iu) + \conj{2}\,(-1)
= 2 - 2 = 0 . \checkmark
\]
Normalizando ($\norm{v_4}^2 = 2 + 4 = 6$, $\norm{v_1}^2 = 2 + 1 =
3$) se obtiene la matriz \omterm{def:b2:hermitian:adjoint}{unitaria} $U = \bigl(\frac{v_4}{\sqrt6}\
\frac{v_1}{\sqrt3}\bigr)$ con $A = U\operatorname{diag}(4,1)U^\dagger$.
Moraleja: la lectura de Rayleigh es inmediata: sobre la esfera unidad
de $\C^2$, $\langle x, Ax\rangle$ recorre $\intcc{1}{4}$, alcanzado en
los dos \omterm{def:b2:reduction:eigen}{vectores propios}; es el germen $n = 2$ de la teoría de
Courant--Fischer que se construye en el problema de fin de semana.
Comprobemos también que la forma es real fuera de los
\omterm{def:b2:reduction:eigen}{vectores propios}: en $x = (1, \iu)$,
\[
Ax = \bigl(2 + (1-\iu)\iu,\ (1+\iu) + 3\iu\bigr)
= (3 + \iu,\ 1 + 4\iu),
\]
\[
\langle x, Ax\rangle = \conj{1}\,(3+\iu) +
\conj{\iu}\,(1+4\iu)
= (3 + \iu) + (-\iu)(1 + 4\iu) = 3 + \iu - \iu + 4 = 7 \in
\R ,
\]
como garantiza para todo $x$ el mecanismo de la demostración de la~\cref{prop:b2:hermitian:eigenvalues} (la simetría conjugada frente a
$A^\dagger = A$).
\end{example}

\begin{example}[Una matriz unitaria diagonalizada]\label{ex:b2:hermitian:unitary}
$U = \frac{1}{\sqrt2}\begin{pmatrix} 1 & \iu\\ \iu & 1
\end{pmatrix}$ (\omterm{def:b2:hermitian:adjoint}{unitaria} por el~\cref{exo:b2:hermitian:2}). Su
\omterm{def:b2:reduction:charpoly}{polinomio característico} es $\bigl(X - \frac{1}{\sqrt2}\bigr)^2 +
\frac12$, de raíces
\[
\lambda_\pm = \frac{1 \pm \iu}{\sqrt2} = \eu^{\pm\iu\pi/4},
\]
de módulo $1$ como prometía la~\cref{prop:b2:hermitian:eigenvalues}, y
con \omterm{def:b2:reduction:eigen}{vectores propios} ortonormales $\frac{1}{\sqrt2}(1, \pm1)$. Así
pues, $U = V\operatorname{diag}(\eu^{\iu\pi/4},
\eu^{-\iu\pi/4})V^\dagger$: en la base adecuada, $U$ es un par de
rotaciones planas de ángulos $\pm\frac\pi4$; una matriz de rotación
real no tiene \omterm{def:b2:reduction:eigen}{vectores propios} reales, pero sobre $\C$ se escinde en
dos escalares de módulo $1$. Moraleja: los \omterm{def:b2:reduction:eigen}{espectros} \omterm{def:b2:hermitian:adjoint}{hermíticos} viven
en la recta real y los \omterm{def:b2:hermitian:adjoint}{unitarios} en la circunferencia unidad; ambos
son sombras de la misma normalidad, y la transformada de Cayley del~\cref{exo:b2:hermitian:6} lleva una imagen a la otra.
\end{example}

\begin{example}[La transformada de Cayley, calculada]\label{ex:b2:hermitian:cayleyrun}
Ejecutemos el~\cref{exo:b2:hermitian:6} sobre $H = \begin{pmatrix} 0
& 1\\ 1 & 0\end{pmatrix}$ (\omterm{def:b2:hermitian:adjoint}{hermítica}, de \omterm{def:b2:reduction:eigen}{espectro} $\{1, -1\}$ y
\omterm{def:b2:reduction:eigen}{vectores propios} ortonormales $\frac{1}{\sqrt2}(1, \pm1)$). En la base
de \omterm{def:b2:reduction:eigen}{vectores propios} todo es escalar: la transformada $\lambda \mapsto
\frac{\lambda - \iu}{\lambda + \iu}$ manda
\[
1 \longmapsto \frac{1 - \iu}{1 + \iu} = -\iu,
\qquad
-1 \longmapsto \frac{-1 - \iu}{-1 + \iu} = \iu ,
\]
(multiplíquese por el conjugado del denominador), de modo que $U = (H
- \iu I)(H + \iu I)^{-1}$ es la matriz \omterm{def:b2:hermitian:adjoint}{unitaria} de
\omterm{def:b2:reduction:eigen}{valores propios} $\mp\iu$ sobre esos mismos \omterm{def:b2:reduction:eigen}{vectores propios}:
\[
U = \frac{1}{2}\begin{pmatrix} 1 & 1\\ 1 & -1\end{pmatrix}
\begin{pmatrix} -\iu & 0\\ 0 & \iu\end{pmatrix}
\begin{pmatrix} 1 & 1\\ 1 & -1\end{pmatrix}
= \begin{pmatrix} 0 & -\iu\\ -\iu & 0\end{pmatrix} .
\]
Comprobación: $U^\dagger U = I$ y $1 \notin \operatorname{Sp}U =
\{\pm\iu\}$, como promete la teoría. Moraleja: la recta real se
aplica sobre la circunferencia unidad menos el punto $1$;
\omterm{def:b2:reduction:eigen}{valor propio} a \omterm{def:b2:reduction:eigen}{valor propio}, la transformada de Cayley es la
aplicación de Möbius $\frac{\lambda-\iu}{\lambda+\iu}$, y las
matrices no hacen más que seguir a su \omterm{def:b2:reduction:eigen}{espectro}.
\end{example}

\begin{remark}[Errores frecuentes]
\emph{(i) Dónde cae la barra:} este libro conjuga la \emph{primera}
entrada, así que las coordenadas son $\langle e_i, x\rangle$ y
$\langle\lambda x, y\rangle = \conj\lambda\langle x, y\rangle$;
muchos textos conjugan en cambio la segunda entrada: tradúzcase antes
de comparar fórmulas, o los signos de $\iu$ se torcerán en silencio.
\emph{(ii) La \omterm{def:b2:quadratic:def}{polarización} compleja es más fuerte:} sobre $\C$, si
$\langle x, u(x)\rangle = 0$ para \emph{todo} $x$, entonces $u = 0$
(desarróllese en $x + y$ y en $x + \iu y$: se anulan tanto la parte
real como la imaginaria de $\langle x, u(y)\rangle$); sobre $\R$ esto
falla: la rotación de ángulo $\frac\pi2$ cumple $\langle x,
u(x)\rangle = 0$ en todas partes. En consecuencia, y solo sobre $\C$,
``$\langle x, u(x)\rangle \in \R$ para todo $x$'' ya fuerza que $u$
sea \omterm{def:b2:hermitian:adjoint}{hermítico}. \emph{(iii) Normal real no significa
\omterm{def:b2:reduction:diag}{diagonalizable}:} la matriz del~\cref{ex:b2:hermitian:skewexample} es normal pero no tiene ningún
\omterm{def:b2:reduction:eigen}{valor propio} real; la diagonalización \omterm{def:b2:hermitian:adjoint}{unitaria} es un teorema
\emph{sobre} $\C$, y sobre $\R$ solo se obtienen reducciones por
bloques. \emph{(iv) Comprobar la unitariedad:} $U^\dagger U = I$
significa que las \emph{columnas} son ortonormales para el producto
\omterm{def:b2:hermitian:adjoint}{hermítico}; probar con $UU^{\mathsf T}$, u olvidar la conjugación en
los productos de columnas, son las dos maneras clásicas de certificar
una matriz equivocada.
\end{remark}

\begin{example}[Las isometrías son exactamente las unitarias]\label{ex:b2:hermitian:isometry}
Conservar la \omterm{def:b2:nvs:norm}{norma} parece más débil que ser \omterm{def:b2:hermitian:adjoint}{unitario}, pero sobre
$\C$ no lo es: si $\norm{u(x)} = \norm x$ para todo $x$, entonces
$u^*u = \mathrm{id}$. En efecto, $v = u^*u - \mathrm{id}$ es
\omterm{def:b2:hermitian:adjoint}{hermítico} y cumple $\langle x, v(x)\rangle = \norm{u(x)}^2 - \norm
x^2 = 0$ para todo $x$; por \omterm{def:b2:quadratic:def}{polarización} compleja (el error frecuente
(ii) de arriba), una aplicación cuya ``diagonal'' se anula
idénticamente es nula: $v = 0$. En concreto, la \omterm{def:b2:quadratic:def}{polarización} se lee
\[
0 = \langle x + y, v(x+y)\rangle
= \langle x, v(y)\rangle + \langle y, v(x)\rangle,
\qquad
0 = \langle x + \iu y, v(x + \iu y)\rangle
= \iu\langle x, v(y)\rangle - \iu\langle y, v(x)\rangle ,
\]
y las dos líneas juntas fuerzan $\langle x, v(y)\rangle = 0$ para
todos $x, y$. Moraleja: por eso ``\omterm{def:b2:hermitian:adjoint}{unitario}'' puede comprobarse
midiendo solo longitudes; una rigidez que el capítulo de Fourier
explotará, donde conservar la energía $\norm f_2$ (Parseval) equivale
a conservar todos los productos escalares de coeficientes.
\end{example}

\begin{remark}[Perspectivas dentro de este volumen]
La maquinaria \omterm{def:b2:hermitian:adjoint}{hermítica} construida aquí se consume casi de
inmediato. El capítulo de Fourier \emph{es} geometría \omterm{def:b2:hermitian:adjoint}{hermítica} en
dimensión infinita: las exponenciales $(e_n)$ forman una familia
ortonormal para $\langle f, g\rangle = \frac{1}{2\pi}\int\conj fg$,
la desigualdad de Bessel es la estimación de proyección del~\cref{thm:b2:hermitian:cs} de este capítulo, y Parseval es su
igualdad en el límite. La transformada de Fourier finita
(\cref{exo:b2:hermitian:10}) reaparece siempre que hay que
diagonalizar una convolución. Y el problema de fin de semana de este
capítulo ---Courant--Fischer, Weyl, entrelazamiento--- suministra la
estabilidad de los \omterm{def:b2:reduction:eigen}{valores propios} que invoca el capítulo de
ecuaciones diferenciales cuando afirma que las pequeñas
perturbaciones de un sistema mueven poco sus frecuencias. Hacia
atrás, todo lo de aquí es el espejo complejo del capítulo de \omterm{def:b2:quadratic:def}{formas
cuadráticas}: manténganse los dos diccionarios uno junto al otro
($A^{\mathsf T} \leftrightarrow A^\dagger$, ortogonal
$\leftrightarrow$ \omterm{def:b2:hermitian:adjoint}{unitario}, Rayleigh real en ambos).
\end{remark}

\begin{remark}[Endomorfismos normales]
Sobre $\C$, el enunciado definitivo es: $u$ es \emph{unitariamente
\omterm{def:b2:reduction:diag}{diagonalizable} si y solo si es normal} ($u^*u = uu^*$), lo que
cubre de una vez las aplicaciones \omterm{def:b2:hermitian:adjoint}{hermíticas}, \omterm{def:b2:hermitian:adjoint}{unitarias} y
antihermíticas. La demostración es un agradable refuerzo del
argumento anterior (\cref{exo:b2:hermitian:8}). Sobre $\R$, en
cambio, la normalidad solo da diagonalización por bloques (bloques de
rotación): la geometría compleja es genuinamente más sencilla.
\end{remark}

\begin{remark}[Dónde se usa]
La teoría espectral \omterm{def:b2:hermitian:adjoint}{hermítica} es la matemática de la mecánica
cuántica: los observables se modelan mediante operadores
\omterm{def:b2:hermitian:adjoint}{hermíticos} (\omterm{def:b2:reduction:eigen}{espectro} real $=$ valores medibles) y la evolución
temporal mediante operadores \omterm{def:b2:hermitian:adjoint}{unitarios} (conservación de la \omterm{def:b2:nvs:norm}{norma}
$=$ conservación de la probabilidad). Dentro de este libro, el
capítulo de Fourier se apoya en la ortonormalidad de las
exponenciales ---un enunciado sobre el producto escalar
hermítico--- y la diagonalización de las matrices circulantes
(\cref{exo:b2:hermitian:10}) es la transformada de Fourier finita. El
problema de fin de semana desarrolla el cálculo variacional de los
\omterm{def:b2:reduction:eigen}{valores propios} (Courant--Fischer, Weyl, entrelazamiento), el pan de
cada día del análisis numérico y de la física matemática; el volumen
del tercer año lo extiende a los operadores \omterm{def:b2:metric:compact}{compactos}
autoadjuntos sobre espacios de Hilbert.
\end{remark}

\section{Ejercicios}

\begin{exercise}[$\star$]\label{exo:b2:hermitian:1}
En $\C^2$: calcula $\langle x, y\rangle$, $\norm x$ y $\norm y$ para
$x = (1, \iu)$, $y = (\iu, 1)$; ¿son ortogonales? Da una base
ortonormal que contenga $\frac{x}{\norm x}$.
\end{exercise}

\begin{exercise}[$\star$]\label{exo:b2:hermitian:2}
¿Cuáles son \omterm{def:b2:hermitian:adjoint}{hermíticas}? ¿\omterm{def:b2:hermitian:adjoint}{unitarias}? ¿normales?
\[
\begin{pmatrix} 1 & \iu\\ -\iu & 2 \end{pmatrix},
\qquad
\frac{1}{\sqrt2}\begin{pmatrix} 1 & \iu\\ \iu & 1\end{pmatrix},
\qquad
\begin{pmatrix} 0 & 1\\ 0 & 0 \end{pmatrix}.
\]
\end{exercise}

\begin{exercise}[$\star$]\label{exo:b2:hermitian:3}
Demuestra que toda matriz $A \in \mathcal{M}_n(\C)$ se escribe de
manera única $A = H + \iu K$ con $H, K$ \omterm{def:b2:hermitian:adjoint}{hermíticas} \emph{(las
``partes real e imaginaria'' $H = \frac{A + A^\dagger}{2}$, $K =
\frac{A - A^\dagger}{2\iu}$)}, y que $A$ es normal si y solo si $H$ y
$K$ conmutan.
\end{exercise}

\begin{exercise}[$\star\star$]\label{exo:b2:hermitian:4}
Diagonaliza en una base ortonormal $A = \begin{pmatrix} 2 & \iu\\
-\iu & 2\end{pmatrix}$ y calcula $A^k$ para $k \in \N$.
\end{exercise}

\begin{exercise}[$\star\star$]\label{exo:b2:hermitian:5}
Demuestra que $U(n)$ es \omterm{def:b2:metric:compact}{compacto} y que la aplicación
\omterm{def:b2:reduction:eigen}{valor propio} es sobreyectiva: todo $\lambda$ de módulo $1$ aparece
para alguna matriz \omterm{def:b2:hermitian:adjoint}{unitaria}. Demuestra que $\det U \in \mathbb{U}$
(la circunferencia unidad) para $U \in U(n)$.
\end{exercise}

\begin{exercise}[$\star\star$]\label{exo:b2:hermitian:6}
(Transformada de Cayley) Sea $H$ \omterm{def:b2:hermitian:adjoint}{hermítica}. Demuestra que $H + \iu I$
es invertible y que $U = (H - \iu I)(H + \iu I)^{-1}$ es \omterm{def:b2:hermitian:adjoint}{unitaria},
con $1 \notin \operatorname{Sp}(U)$. \emph{(Trabaja espectralmente:
en una base de \omterm{def:b2:reduction:eigen}{vectores propios} de $H$ todo es escalar.)}
\end{exercise}

\begin{exercise}[$\star\star$]\label{exo:b2:hermitian:7}
Para $A$ \omterm{def:b2:hermitian:adjoint}{hermítica} definida positiva ($\langle x, Ax\rangle > 0$ para
$x \neq 0$), demuestra que $\operatorname{Sp}(A) \subseteq
\intoo{0}{\infty}$, que $A = B^2$ para cierta $B$ \omterm{def:b2:hermitian:adjoint}{hermítica} definida
positiva y que $\det A > 0$.
\end{exercise}

\begin{exercise}[$\star\star\star$]\label{exo:b2:hermitian:8}
(Teorema espectral para endomorfismos normales) Sea $u$ normal sobre
un \omterm{def:b2:hermitian:def}{espacio hermítico}.
\begin{enumerate}
  \item Demuestra que $\norm{u(x)} = \norm{u^*(x)}$ para todo $x$ y
        deduce $\ker(u - \lambda) = \ker(u^* - \conj\lambda)$.
  \item Demuestra que los \omterm{def:b2:reduction:eigen}{subespacios propios} de $u$ para
        \omterm{def:b2:reduction:eigen}{valores propios} distintos son ortogonales, y que el
        complemento ortogonal de un \omterm{def:b2:reduction:eigen}{subespacio propio} es estable por
        $u$.
  \item Concluye por inducción que $u$ es unitariamente
        \omterm{def:b2:reduction:diag}{diagonalizable}; y recíprocamente.
\end{enumerate}
\end{exercise}

\begin{exercise}[$\star$]\label{exo:b2:hermitian:9}
Un endomorfismo es \emph{antihermítico} cuando $u^* = -u$. Demuestra
que sus \omterm{def:b2:reduction:eigen}{valores propios} son imaginarios puros, que $u \mapsto \iu u$
es una biyección de los \omterm{def:b2:hermitian:adjoint}{endomorfismos hermíticos} sobre los
antihermíticos, y que los endomorfismos antihermíticos son
unitariamente \omterm{def:b2:reduction:diag}{diagonalizables} (\cref{exo:b2:hermitian:8}).
\end{exercise}

\begin{exercise}[$\star\star$]\label{exo:b2:hermitian:10}
(La transformada de Fourier finita) Sea $S$ el desplazamiento
\omterm{def:b2:structures:generated}{cíclico} de $\C^n$: $S(x_0, x_1, \dots, x_{n-1}) = (x_{n-1}, x_0,
\dots, x_{n-2})$, y sea $\omega = \eu^{2\iu\pi/n}$.
\begin{enumerate}
  \item Prueba que $S$ es \omterm{def:b2:hermitian:adjoint}{unitario} y que los vectores $f_k =
        \frac{1}{\sqrt n}\bigl(1, \omega^k, \omega^{2k}, \dots,
        \omega^{(n-1)k}\bigr)$, $0 \leq k < n$, forman una base
        ortonormal de \omterm{def:b2:reduction:eigen}{vectores propios}: $Sf_k = \omega^{-k} f_k$.
  \item Deduce que toda matriz \emph{circulante} $C =
        \sum_{j=0}^{n-1} c_jS^j$ es normal, se diagonaliza en la
        misma base y tiene \omterm{def:b2:reduction:eigen}{valores propios} $\widehat c(k) = \sum_j
        c_j\,\omega^{-jk}$.
\end{enumerate}
\end{exercise}

\begin{exercise}[$\star\star$]\label{exo:b2:hermitian:11}
Sea $P$ un endomorfismo idempotente ($P^2 = P$) de un \omterm{def:b2:hermitian:def}{espacio
hermítico}. Demuestra que $P$ es la proyección \emph{ortogonal} sobre
$\operatorname{im} P$ si y solo si $P^* = P$. Da la matriz de la
proyección ortogonal sobre $\C v$ ($\norm v = 1$), y sobre un
subespacio de base ortonormal $(v_1, \dots, v_k)$.
\end{exercise}

\begin{exercise}[$\star\star\star$]\label{exo:b2:hermitian:12}
(Proyectores espectrales por interpolación) Sea $A$ \omterm{def:b2:hermitian:adjoint}{hermítica} con
\omterm{def:b2:reduction:eigen}{valores propios} \emph{distintos} $\lambda_1, \dots, \lambda_p$ y
descomposición en \omterm{def:b2:reduction:eigen}{subespacios propios} $E = \bigoplus_i E_i$.
Definimos los polinomios de Lagrange $L_i(X) = \prod_{j\neq i}\frac{X
- \lambda_j}{\lambda_i - \lambda_j}$. Demuestra que $P_i = L_i(A)$ es
la proyección ortogonal sobre $E_i$, que $P_iP_j = 0$ para $i \neq
j$, que $\sum_i P_i = I$ y que $A = \sum_i \lambda_iP_i$ (la
\emph{descomposición espectral}); expresa $f(A)$ en términos de los
$P_i$ para un polinomio $f$ cualquiera.
\end{exercise}

\section{Problema: Courant--Fischer, Weyl y el cálculo de valores
propios}

\begin{problem}\label{pb:b2:hermitian:1}
Los \omterm{def:b2:reduction:eigen}{valores propios} de una matriz \omterm{def:b2:hermitian:adjoint}{hermítica} no son solo raíces de un
polinomio: son soluciones de \emph{problemas de optimización}. Ese
punto de vista variacional ---los \omterm{pb:b2:hermitian:1}{cocientes de Rayleigh} y el
\emph{teorema mín-máx de Courant--Fischer}\index{teorema de
Courant--Fischer}--- hace que los \omterm{def:b2:reduction:eigen}{valores propios} sean comparables,
estables y calculables, y este problema recoge sus cosechas clásicas:
las \emph{desigualdades de perturbación de Weyl}\index{desigualdad de
Weyl}, el \emph{entrelazamiento de Cauchy}, las desigualdades de
trazas de Schur y de Ky Fan, la monotonía de la raíz cuadrada
matricial y el \omterm{def:b2:reduction:eigen}{espectro} del laplaciano discreto. En todo el
problema, $A, B, E$ son \omterm{def:b2:hermitian:adjoint}{hermíticas} sobre $E = \C^n$ con
\omterm{def:b2:reduction:eigen}{valores propios} listados en orden decreciente $\lambda_1(A) \geq
\dots \geq \lambda_n(A)$, y $R_A(x) = \frac{\langle x, Ax\rangle}
{\langle x, x\rangle}$ para $x \neq 0$ es el \emph{cociente de
Rayleigh}\index{cociente de Rayleigh}.

\medskip
\noindent\textbf{Parte I --- \omterm{pb:b2:hermitian:1}{Cocientes de Rayleigh} y mín-máx.}
Fijemos una base ortonormal de \omterm{def:b2:reduction:eigen}{vectores propios} $(e_1, \dots, e_n)$,
$Ae_i = \lambda_ie_i$.
\begin{enumerate}
  \item Prueba que $R_A(x)$ es real y que
        \[
        \lambda_n \leq R_A(x) \leq \lambda_1
        \qquad (x \neq 0),
        \]
        alcanzándose ambas cotas: $\lambda_1 = \max R_A$ y
        $\lambda_n = \min R_A$.
  \item Prueba que los puntos críticos de $R_A$ son exactamente los
        \omterm{def:b2:reduction:eigen}{vectores propios} de $A$ \emph{(desarrolla $t \mapsto R_A(x +
        tv)$ en $t = 0$ para $v$ arbitrario y sustituye después $v$
        por $\iu v$)}.
  \item Sean $V_k = \operatorname{Vect}(e_1, \dots, e_k)$ y $W_k =
        \operatorname{Vect}(e_k, \dots, e_n)$. Prueba que
        \[
        \min_{x \in V_k\setminus\{0\}} R_A(x) = \lambda_k
        = \max_{x \in W_k\setminus\{0\}} R_A(x) .
        \]
  \item Demuestra el \emph{teorema de Courant--Fischer}: para $1
        \leq k \leq n$,
        \[
        \lambda_k = \max_{\dim V = k}\;\min_{x \in
        V\setminus\{0\}} R_A(x)
        = \min_{\dim W = n-k+1}\;\max_{x \in
        W\setminus\{0\}} R_A(x)
        \]
        \emph{(para todo $V$ de dimensión $k$: $V \cap W_k \neq
        \{0\}$ por Grassmann, luego $\min_V R_A \leq \lambda_k$; la
        pregunta 3 muestra que la cota se alcanza)}.
  \item (Monotonía) Escribimos $A \leq B$ cuando $B - A$ es
        semidefinida positiva. Deduce de la pregunta 4 que $A \leq
        B$ implica $\lambda_k(A) \leq \lambda_k(B)$ para todo $k$.
\end{enumerate}

\medskip
\noindent\textbf{Parte II --- Las desigualdades de Weyl.}
\begin{enumerate}[resume]
  \item Prueba que dos subespacios $V, W \subseteq \C^n$ con $\dim V
        + \dim W > n$ se cortan de manera no trivial, y generaliza:
        $\dim(V_1 \cap V_2 \cap V_3) \geq \dim V_1 + \dim V_2 + \dim
        V_3 - 2n$.
  \item Demuestra la \emph{desigualdad de Weyl}: para $i + j - 1
        \leq n$,
        \[
        \lambda_{i+j-1}(A + B) \leq \lambda_i(A) +
        \lambda_j(B)
        \]
        \emph{(corta los subespacios $W_i(A)$, $W_j(B)$ y
        $V_{i+j-1}(A+B)$ de la pregunta 3 y cuenta dimensiones)}.
  \item Define $\vertiii{E}_2 = \max_{\norm x = 1}\norm{Ex}$ y
        prueba que $\vertiii E_2 = \max_k\abs{\lambda_k(E)}$ para
        $E$ \omterm{def:b2:hermitian:adjoint}{hermítica}. Deduce el \emph{teorema de perturbación de
        Weyl}:
        \[
        \bigl|\lambda_k(A + E) - \lambda_k(A)\bigr| \leq
        \vertiii{E}_2
        \qquad (1 \leq k \leq n) :
        \]
        cada \omterm{def:b2:reduction:eigen}{valor propio} es una función $1$-lipschitziana de la
        matriz.
  \item (Perturbaciones de rango uno) Sea $P$ \omterm{def:b2:hermitian:adjoint}{hermítica}
        semidefinida positiva de rango $1$. Prueba que
        \[
        \lambda_k(A) \leq \lambda_k(A + P) \leq
        \lambda_{k-1}(A) \qquad (2 \leq k \leq n),
        \]
        junto con $\lambda_1(A) \leq \lambda_1(A+P)$: los nuevos
        \omterm{def:b2:reduction:eigen}{valores propios} \emph{se entrelazan} con los antiguos.
  \item Comprueba numéricamente la pregunta 8: $A = \begin{pmatrix}
        2 & \iu\\ -\iu & 2\end{pmatrix}$
        (\cref{exo:b2:hermitian:4}: \omterm{def:b2:reduction:eigen}{espectro} $\{3, 1\}$) y $E =
        \begin{pmatrix} 0 & 1\\ 1 & 0\end{pmatrix}$ (\omterm{def:b2:reduction:eigen}{espectro}
        $\{1, -1\}$): calcula el \omterm{def:b2:reduction:eigen}{espectro} de $A + E$ y los dos
        miembros de la desigualdad.
\end{enumerate}

\medskip
\noindent\textbf{Parte III --- Entrelazamiento y desigualdades de
trazas.}
\begin{enumerate}[resume]
  \item (Entrelazamiento de Cauchy) Sea $B$ la submatriz principal
        dominante de tamaño $(n-1)\times(n-1)$ de $A$. Demuestra que
        \[
        \lambda_{k+1}(A) \leq \lambda_k(B) \leq \lambda_k(A)
        \qquad (1 \leq k \leq n-1)
        \]
        \emph{(considera $\C^{n-1} \subseteq \C^n$; sobre él, $R_B$
        es la restricción de $R_A$; aplica Courant--Fischer en los
        dos niveles)}.
  \item Itera: para una submatriz principal $B$ de tamaño $n - m$,
        $\lambda_{k+m}(A) \leq \lambda_k(B) \leq \lambda_k(A)$.
  \item (Schur) Sean $d_1 \geq d_2 \geq \dots \geq d_n$ las entradas
        diagonales de $A$, ordenadas. Demuestra, para todo $k$, que
        \[
        \sum_{i=1}^{k} d_i \leq \sum_{i=1}^{k}\lambda_i(A),
        \]
        con igualdad en $k = n$ (la traza) \emph{(las $k$ entradas
        diagonales elegidas forman una submatriz principal $k\times
        k$; acota su traza mediante la pregunta 12)}.
  \item (Ky Fan) Demuestra que
        \[
        \sum_{i=1}^{k}\lambda_i(A)
        = \max\Bigl\{\sum_{i=1}^{k}\langle x_i, Ax_i\rangle
        : (x_1, \dots, x_k) \text{ ortonormal}\Bigr\} .
        \]
  \item Verifica las preguntas 11 y 13 sobre
        \[
        A = \begin{pmatrix} 2 & 1 & 0\\ 1 & 2 & 1\\
        0 & 1 & 2\end{pmatrix}
        \qquad
        \bigl(\operatorname{Sp} = \{2 - \sqrt2,\ 2,\
        2+\sqrt2\}\bigr)
        \]
        frente a su bloque dominante $2\times2$ (de \omterm{def:b2:reduction:eigen}{espectro} $\{1,
        3\}$) y a su diagonal.
\end{enumerate}

\medskip
\noindent\textbf{Parte IV --- El orden de Loewner.} $A \leq B$ sigue
significando que $B - A$ es semidefinida positiva; todas las matrices
de esta parte son \omterm{def:b2:hermitian:adjoint}{hermíticas}.
\begin{enumerate}[resume]
  \item Prueba que $A \leq B$ implica $a_{ii} \leq b_{ii}$ para todo
        $i$, $\operatorname{tr} A \leq \operatorname{tr} B$ y
        $C^\dagger AC \leq C^\dagger BC$ para \emph{toda} matriz
        compleja $C$.
  \item Prueba que elevar al cuadrado \emph{no} es monótono: para
        \[
        A = \begin{pmatrix} 1 & 0\\ 0 & 0\end{pmatrix},
        \qquad
        B = \begin{pmatrix} 2 & 1\\ 1 & 1\end{pmatrix},
        \]
        comprueba que $0 \leq A \leq B$ pero $A^2 \not\leq B^2$.
  \item Demuestra que la raíz cuadrada \emph{sí} es monótona: $0
        \leq A \leq B$ implica $\sqrt A \leq \sqrt B$ \emph{(sea
        $\mu$ un \omterm{def:b2:reduction:eigen}{valor propio} de $\sqrt B - \sqrt A$ con
        \omterm{def:b2:reduction:eigen}{vector propio} \omterm{def:b2:hermitian:adjoint}{unitario} $v$; calcula $\langle v, (B -
        A)v\rangle = \mu\bigl(\langle v, \sqrt B\,v\rangle + \langle
        v, \sqrt A\,v\rangle\bigr)$ y discute)}.
  \item Demuestra que la inversión es antítona sobre las matrices
        definidas positivas: $0 < A \leq B$ implica $B^{-1} \leq
        A^{-1}$ \emph{(hágase la \omterm{def:b2:quadratic:def}{congruencia} por $A^{-1/2}$ para
        reducirse a $I \leq M \Rightarrow M^{-1} \leq I$, que es
        escalar en una base espectral)}.
  \item Sean $A, B$ definidas positivas. Prueba que los
        \omterm{def:b2:reduction:eigen}{valores propios} de $AB$ (¡no \omterm{def:b2:hermitian:adjoint}{hermítica} en general!) son
        reales y positivos, y que
        \[
        \lambda_{\max}(AB) \leq
        \lambda_{\max}(A)\,\lambda_{\max}(B)
        \]
        \emph{(conjuga por $\sqrt A$: $AB \sim \sqrt A\,B\sqrt
        A$)}.
\end{enumerate}

\medskip
\noindent\textbf{Parte V --- El laplaciano discreto, resuelto.} Sea
$T_n$ la matriz tridiagonal $n \times n$ con $2$ en la diagonal y
$-1$ en las dos diagonales adyacentes.
\begin{enumerate}[resume]
  \item Con $\theta_k = \frac{k\pi}{n+1}$, comprueba que los
        vectores $v_k = \bigl(\sin(j\theta_k)\bigr)_{1\leq j\leq n}$
        cumplen $T_nv_k = (2 - 2\cos\theta_k)\,v_k$ \emph{(identidad
        de producto a suma; compruébense las filas de los extremos,
        $j = 1, n$)}. Concluye que
        \[
        \operatorname{Sp}(T_n) = \Bigl\{4\sin^2
        \frac{k\pi}{2(n+1)} : 1 \leq k \leq n\Bigr\},
        \]
        todos simples y positivos: $T_n$ es definida positiva.
  \item (Un potencial) Para una matriz diagonal real $D =
        \operatorname{diag}(d_1, \dots, d_n)$, encierra el
        \omterm{def:b2:reduction:eigen}{espectro}: para todo $k$,
        \[
        \lambda_k(T_n) + \min_i d_i \;\leq\;
        \lambda_k(T_n + D) \;\leq\; \lambda_k(T_n) + \max_i
        d_i .
        \]
  \item Comprueba explícitamente el entrelazamiento de Cauchy entre
        $T_3$ y $T_2$ (de \omterm{def:b2:reduction:eigen}{espectros} $\{2 \pm \sqrt2, 2\}$ y $\{1,
        3\}$) e interprétalo: $T_2$ es $T_3$ tras suprimir un
        extremo del camino.
  \item Prueba que los \omterm{def:b2:reduction:eigen}{valores propios} extremos cumplen, cuando $n
        \to \infty$,
        \[
        \lambda_{\min}(T_n) = 4\sin^2\frac{\pi}{2(n+1)}
        \sim \frac{\pi^2}{(n+1)^2},
        \qquad
        \lambda_{\max}(T_n) \to 4 ,
        \]
        de modo que el número de condición $\kappa_n =
        \lambda_{\max}/\lambda_{\min}$ crece como
        $\frac{4(n+1)^2}{\pi^2}$: discretizar una derivada segunda
        sobre una malla cada vez más fina está intrínsecamente mal
        condicionado.
  \item Síntesis. Una frase para cada punto: (i) por qué es la
        caracterización variacional, y no el \omterm{def:b2:reduction:charpoly}{polinomio
        característico}, lo que hace estables a los
        \omterm{def:b2:reduction:eigen}{valores propios} (preguntas 8--9); (ii) qué preguntas usaron
        solo $\lambda_1 = \max R_A$ y cuáles necesitaron el mín-máx
        \omterm{def:b2:metric:complete}{completo}; (iii) qué añade a la historia el orden de Loewner;
        (iv) dónde reaparecen estas herramientas (análisis numérico
        de las matrices de rigidez de la pregunta 24; teoría
        cuántica de perturbaciones; y, en el volumen del tercer año,
        el principio mín-máx para operadores \omterm{def:b2:metric:compact}{compactos}
        autoadjuntos).
\end{enumerate}
\end{problem}
